% Théorème de l'énergie cinétique pour une luge qui descend une pente avec frottements (L'essentiel, chapitre 15)
\begin{tikzpicture}[x=1cm,y=1cm,
  force/.style={line width=1.2pt,-{Stealth[length=2.4mm]}}]
\fill[otMuted!20] (0,0) -- (7.4,0) -- (0,3.2) -- cycle;
\draw[otInk!70,line width=0.6pt] (0,3.2) -- (7.4,0) -- (0,0);
\begin{scope}[shift={(3.0,1.903)},rotate=-23.4]
  \filldraw[fill=otOrange!60,draw=otOrange!80!black,line width=0.5pt,rounded corners=0.8mm] (-0.45,0.02) rectangle (0.45,0.38);
  \coordinate (G) at (0,0.2);
\end{scope}
\fill[otInk] (G) circle (0.045);
\draw[force,otRed!80!black] (G) -- ++(0,-1.3) node[pcLegende,anchor=north,text=otRed!80!black] {$\vec{P}$};
\draw[force,otNavy] (G) -- ++(66.6:1.15) node[pcLegende,anchor=south,text=otNavy] {$\vec{R}$};
\draw[force,otOrange!85!black] (G) -- ++(156.6:1.0) node[pcLegende,anchor=south,text=otOrange!85!black] {$\vec{f}$};
\fill[otInk] (1.2,2.684) circle (0.05); \node[pcLegendeG,anchor=south west] at (1.2,2.684) {A};
\fill[otInk] (6.0,0.605) circle (0.05); \node[pcLegendeG,anchor=south west] at (6.0,0.605) {B};
\draw[otGreenDark,line width=0.9pt,-{Stealth[length=2mm]}] (6.0,0.85) -- ++(-23.4:1.2) node[pcLegende,anchor=south,text=otGreenDark] {$\vec{v}_\mathrm{B}$};
\node[pcLegende,anchor=west,align=left] at (7.6,1.9) {$W(\vec{P})>0$ : moteur\\$W(\vec{R})=0$ : $\vec{R}\perp$ déplacement\\$W(\vec{f})<0$ : résistant};
\end{tikzpicture}
\par\smallskip\legende{Entre A et B : $E_\mathrm{c}(\mathrm{B})-E_\mathrm{c}(\mathrm{A})=W_{\mathrm{AB}}(\vec{P})+W_{\mathrm{AB}}(\vec{R})+W_{\mathrm{AB}}(\vec{f})$.}
